\subsection{Creating Relations} \label{s:cr}
Armed with the \md{} from \S\ref{s:mxe} we move to the next step of using this \md{} to create relations among objects. For this discussion we use the example of an auto insurance company servicing a customer claim after a vehicle collision. We have three components to the claim: the claim form, digital photographs taken by the appraiser of the damage sustained, and a scan of the police report. These rich data types are good candidates for an \os{} as all are \ud{}.

Traditional \apps{} create a \db{} to collect transactions and create a schema to associate\db{} tables with one another. However, mobile \apps{} may not wish to require a local \db{}, instead choosing a remote \db{} for this function. In \S\ref{s:uddb} we state that \oss{} are a \db{} for \ud{}. Here RDOS acts as the \db{} for both the scalar data (in the form of key-value \md{} tags) and the data objects (blobs), which consist of any arbitrary file type. To use RDOS in our insurance example we begin with the claim form where a Claim ID (\texttt{CID}) is assigned (\texttt{CID=1234}) and ingested into RDOS. Subsequently the photos from the appraiser and the scan of the police report are uploaded. For all objects there is a \md{} tag set where \texttt{CID=1234}, providing a common key amongst the objects. At ingest, the client \app{} can choose to add a second \md{} field indicating that the images are related to the claim form. As objects are identified by URI, this translates to: \texttt{RelTo=URI\{Object1\}}. Through this structure we have a means of relating the full set of objects which are all associated with the same insurance claim.

This example depends upon the \app{} explicitly creating the relation between the objects by setting the \texttt{RelTo} \md{} tag, which is the preferred method for new object ingest. However, the same mechanism described in \S\ref{s:mxe} for adding \md{} to objects already stored can be applied here. The final stage of a MGM pipeline can be used to assign the relationship tag to objects in the same manner as all other \md{}. The difference in this final MGM script is that it must persist all \md{} tags found when enumerating over a set of existing objects. The system administrator can then define how these objects are to be related via a GUI, which in turn creates the needed MGM script.
